% ---------------- APPENDIX ----------------
\clearpage
\phantomsection
\addcontentsline{toc}{chapter}{Appendix}
\begin{center}
{\Large\bfseries Appendix}
\end{center}

\vspace{0.5cm}

% The reference report uses manual "A / A1. / A2." headings rather than
% LaTeX's \appendix numbering. Keeping the same convention here.

\phantomsection
\section*{A. Project Management and Development Screenshots}
\addcontentsline{toc}{section}{A. Project Management and Development Screenshots}

\phantomsection
\subsection*{A1. Asana Kanban Board}
\addcontentsline{toc}{subsection}{A1. Asana Kanban Board}
\begin{figure}[H]
\centering
\includegraphics[width=\textwidth]{asana_board}
\caption{Asana Kanban board used to manage the Agile sprints}
\label{fig:asana-board}
\end{figure}


\phantomsection
\subsection*{A2. Asana Task Detail}
\addcontentsline{toc}{subsection}{A2. Asana Task Detail}
\begin{figure}[H]
\centering
\includegraphics[height=0.72\textheight,width=\textwidth,keepaspectratio]{asana_task}
\caption{Task detail view showing assignee, subtasks and activity history}
\label{fig:asana-task}
\end{figure}


\phantomsection
\subsection*{A3. Emulator Snapshot}
\addcontentsline{toc}{subsection}{A3. Emulator Snapshot}
\begin{figure}[H]
\centering
\includegraphics[width=0.62\textwidth]{ram_emulator}
\caption{Emulator snapshot of the rendered game frame}
\label{fig:ram-emulator}
\end{figure}


\phantomsection
\subsection*{A4. Raw WRAM Tilemap}
\addcontentsline{toc}{subsection}{A4. Raw WRAM Tilemap}
\begin{figure}[H]
\centering
\includegraphics[width=0.95\textwidth]{ram_hex}
\caption{Raw WRAM tilemap read from emulator memory, shown as hex values}
\label{fig:ram-hex}
\end{figure}


\phantomsection
\subsection*{A5. Quantized Input Grid}
\addcontentsline{toc}{subsection}{A5. Quantized Input Grid}
\begin{figure}[H]
\centering
\includegraphics[width=0.95\textwidth]{ram_quantized}
\caption{The $16 \times 13$ quantized grid supplied to the network, with Mario
marked M and an enemy marked E}
\label{fig:ram-quantized}
\end{figure}


\phantomsection
\subsection*{A6. Training Visualisation Web UI}
\addcontentsline{toc}{subsection}{A6. Training Visualisation Web UI}
\begin{figure}[H]
\centering
\includegraphics[width=\textwidth]{webui_training}
\caption{Web interface used to configure runs and visualise training}
\label{fig:webui-training}
\end{figure}


\clearpage
\phantomsection
\section*{B. Assistance for Running the Project}
\addcontentsline{toc}{section}{B. Assistance for Running the Project}

The project utilizes docker Compose to maintain a consistent execution
environment across all operating systems. This ensures that dependencies and
configurations remain isolated and reproducible.

\phantomsection
\subsection*{B1. Prerequisites}
\addcontentsline{toc}{subsection}{B1. Prerequisites}

Ensure that the following software is installed on your host machine:

\begin{itemize}
    \item Git for version control.
    \item docker Desktop (or docker Engine with docker Compose) for
    containerization.
\end{itemize}

\phantomsection
\subsection*{B2. Project Initialization}
\addcontentsline{toc}{subsection}{B2. Project Initialization}

\noindent
\textbf{1. Clone the Repository}

\begin{lstlisting}
git clone https://gitlab.almondsai.com/khce/neat-for-mario.git
cd neat-for-mario
\end{lstlisting}

\noindent
\textbf{2. Add Game ROMs}

\noindent
Due to copyright and legal constraints, the original Nintendo game ROMs cannot
be tracked in version control.

\begin{enumerate}
    \item Create a folder named \texttt{roms/} in the project root directory.
    \item Place the legally acquired Super Mario Bros.\ ROM inside this folder
    before attempting to run the project.
\end{enumerate}

\noindent
\textbf{3. Environment Configuration (Linux Users Only)}

\noindent
If running natively on Linux (Ubuntu, Arch, etc.), generate a user environment
file before building the container to ensure file permissions are correctly
mapped from the docker container to the host filesystem:

\begin{lstlisting}
echo -e "HOST_UID=$(id -u)\nHOST_GID=$(id -g)" > .env
\end{lstlisting}

\noindent
\textbf{4. Build the docker Image}

\begin{lstlisting}
docker compose build
\end{lstlisting}

\phantomsection
\subsection*{B3. Daily Development Workflow}
\addcontentsline{toc}{subsection}{B3. Daily Development Workflow}

Code editing should be done on the host machine using standard IDEs (VS Code,
Neovim), while code execution takes place entirely inside the docker container.

\noindent
\textbf{1. Accessing the Container Shell}

\noindent
Spin up the container and access its interactive terminal:

\begin{lstlisting}
docker compose run --rm app
\end{lstlisting}

\noindent
\textbf{Note:} The terminal prompt will change to indicate execution is now
happening inside the docker environment.

\noindent
\textbf{2. Syncing Dependencies and Execution}

\noindent
Ensure the python environment is synced with the locked configuration using
\texttt{uv}, then activate the environment and run the training pipeline:

\begin{lstlisting}
# Sync dependencies
uv sync --frozen

# Activate virtual environment and run
source .venv/bin/activate
python main.py
\end{lstlisting}

\phantomsection
\subsection*{B4. Interactive Play Mode}
\addcontentsline{toc}{subsection}{B4. Interactive Play Mode}

To verify the environment or simply play the game manually using keyboard
controls, you must first import the ROM into stable retro, then launch the
interactive script:

\noindent
\textbf{1. Import the ROM}

\begin{lstlisting}
uv run python -m stable_retro.import
\end{lstlisting}

\noindent
\textbf{2. Launch the Game}

\begin{lstlisting}
uv run python -m stable_retro.examples.interactive \
    --game SuperMarioBros-Nes-v0
\end{lstlisting}
